% 发给学生的教材版式：B5、中文、书签、代码可断行。
\usepackage{geometry}
\geometry{b5paper,margin=2.1cm,top=2.2cm,bottom=2.3cm}
\usepackage{fontspec}
\usepackage[AutoFallBack=true]{xeCJK}
% 字体钉死为 Noto 一家（D-042，人于 2026-09-19 拍板），任何机器上排出同一本书。
% 此前是「macOS 用 Times/Menlo/宋体-黑体，别处逐级回退」：同一份源在 Mac 上排 720 页、
% 在 Linux 上排 700 页；Mac 那档的粗体还被映射成 Heiti SC Light，强调比正文又小又细。
% 所以**不再回退**：缺哪款字体就当场报错停机，而不是悄悄换一款排出另一本书。
% 全部 SIL OFL，可随 PDF 分发。安装：Ubuntu `apt install fonts-noto-core fonts-noto-mono
% fonts-noto-cjk fonts-noto-extra`；macOS 从 https://fonts.google.com/noto 或 `brew install --cask
% font-noto-serif font-noto-sans-mono font-noto-serif-cjk-sc font-noto-sans-cjk-sc font-noto-sans-math
% font-noto-sans-symbols-2`。
\newcommand{\dsaRequireFont}[1]{%
  \IfFontExistsTF{#1}{}{%
    \errmessage{缺字体「#1」：学生 PDF 的字体已钉死为 Noto（D-042），不回退。
      请先安装它（见 book/pdf/preamble.tex 开头的安装说明）}}}
\dsaRequireFont{Noto Serif}
\dsaRequireFont{Noto Sans Mono}
\dsaRequireFont{Noto Serif CJK SC}
\dsaRequireFont{Noto Sans CJK SC}
\dsaRequireFont{Noto Sans Math}
\dsaRequireFont{Noto Sans Symbols 2}
\setmainfont{Noto Serif}
\setmonofont{Noto Sans Mono}[Scale=0.82]
\setCJKmainfont{Noto Serif CJK SC}
\setCJKsansfont{Noto Sans CJK SC}
\setCJKmonofont{Noto Sans CJK SC}[Scale=0.88]
% Noto Serif（拉丁）没有正文里用到的这些符号：→ ↔ ⇔ − ∞ ≈ ≤ ≥ ≠ ① ★ ✓ ✗ ∎ ⟹。
% xelatex 不会逐字回退，缺字就是一个空框（钉字体那一版的构建日志里 678 处）。所以把这几段
% 交给 xeCJK，由 Noto Serif CJK SC 排；它也没有的 ∎ ⟹ ✗ 再按固定次序落到 Noto Sans Math、
% Noto Sans Symbols2——回退链是写死的，不随机器变。
\xeCJKDeclareCharClass{CJK}{"2190->"21FF, "2200->"22FF, "2460->"24FF, "2600->"27BF, "27F0->"27FF}
\setCJKfallbackfamilyfont{\CJKrmdefault}{{Noto Sans Math},{Noto Sans Symbols 2}}
\setCJKfallbackfamilyfont{\CJKsfdefault}{{Noto Sans Math},{Noto Sans Symbols 2}}
\setCJKfallbackfamilyfont{\CJKttdefault}{{Noto Sans Math},{Noto Sans Symbols 2}}
% 公式里直接写的汉字（「行」「列」「长度」……共 24 个字、38 处）落在数学字体 Latin Modern Math 上，
% 那里没有汉字——两种旧版式都是空框。xeCJK 的 CJKmath 让公式里的汉字也走 CJK 字体。
% （unicode-math 的 \setmathfont[range=…] 管不到：汉字不在它的数学字符表里，试过无效。）
\xeCJKsetup{CJKspace=true, CJKmath=true}

\usepackage{xcolor}
\PassOptionsToPackage{bookmarks=true,bookmarksnumbered=true,bookmarksopen=true,bookmarksopenlevel=3,colorlinks=true,linkcolor=black,urlcolor=blue!50!black}{hyperref}

\usepackage{fvextra}
\usepackage{fancyvrb}
\DefineVerbatimEnvironment{Highlighting}{Verbatim}{
  breaklines,
  breakanywhere,
  commandchars=\\\{\},
  fontsize=\small,
  frame=single,
  framesep=4pt,
  rulecolor=\color{black!20},
}
\DefineVerbatimEnvironment{verbatim}{Verbatim}{
  breaklines,
  breakanywhere,
  fontsize=\small,
}

\usepackage{longtable}
\usepackage{booktabs}
\usepackage{array}
\usepackage{graphicx}
\usepackage{grffile}
\usepackage{float}
\usepackage{caption}
\captionsetup{font=small,labelfont=bf,skip=6pt}
\DeclareCaptionLabelSeparator{colon}{：}
\captionsetup{labelsep=colon}
\graphicspath{{../}{../../}}
% 插图默认约版心三分之二。pandoc 的 \pandocbounded 只在「超出版心」时缩小，
% 原图更宽时仍会撑满；这里改成硬上限。
\setkeys{Gin}{width=0.62\linewidth,height=0.62\textheight,keepaspectratio}

\usepackage{enumitem}
\setlist{nosep,leftmargin=1.6em}

\pagestyle{plain}
\setlength{\parskip}{0.25em}
\setlength{\parindent}{2em}
\renewcommand{\baselinestretch}{1.12}
% book/ctexbook 默认 flushbottom，标题弹性间距会被拉成整页空白。
\raggedbottom
% pandoc 模板会装 parskip.sty，把段距拉到 6pt plus 2pt；这里压回去。
\setlength{\parskip}{0.2em plus 0.05em minus 0.05em}

% 章首页不要再印一遍巨大的 \chapter 装饰标题：正文里已有「第N章」。
\usepackage{titlesec}
\titleformat{\chapter}[hang]
  {\normalfont\Large\bfseries}
  {}
  {0pt}
  {}
\titleformat{\section}{\normalfont\large\bfseries}{}{0pt}{}
\titleformat{\subsection}{\normalfont\normalsize\bfseries}{}{0pt}{}
\titlespacing*{\chapter}{0pt}{-12pt}{8pt}
\titlespacing*{\section}{0pt}{1.1em}{0.45em}
\titlespacing*{\subsection}{0pt}{0.85em}{0.3em}

\newcommand{\covernote}[1]{\begin{center}\large #1\end{center}}
